% R. Brown's translations of two letters to L. Breen, with his note and comments (cote 134-2): archive cote 134-2, site page 140 (PDF page 141).
% translation of the letter of 17-19 July 1975. Draft compared: ps1.tex, lines 1597-1614.
% AI-transcribed by Claude Opus 5.5 (claude-opus-5-5, adaptive thinking, effort high), Message Batches API, batch brown_full (msgbatch_01Md3i5yK1PxkpJn5GHfEExE), 2026-11-13; not yet checked by eye.
\scanpage{134-2/140}
\mspage{- 13 -}

with boundary.  When in a situation of sheaves (or complexes
of sheaves or n-stacks) which are $P$-constructible on $X$ , where
$P$ is such a fixed stratification, these also can be made
explicit in terms of the envisaged diagram, from which there is
no way to retrieve more than the ``homotopy type''. One thought
that the six operations on sheaves can be translated ad hoc to
this homotopical framework. Finally, if instead of having only
one compact tame space $X$ , one has a tame morphism $f : X \to Y$
then by choosing equisingular stratifications on $X$ and $Y$
adapted to $f$ (the strata of $X$ being notably locally trivial
fibrations on those of $Y$ ...), one should find a ``morphism''
of diagrams of manifolds with boundary expressing $X$ in that
expressing $Y$ (with some natural morphisms essentially equivalent
to fibrations of manifolds with boundaries on the others) in
such a way that the four operations $Rf_*$, $Rf_!$, $Lg^*$, $g^!$ between
$P_X$- and $P_Y$- constructible sheaves on $X$ and $Y$ (on locally
closed sub spaces $X'$ , $Y'$ which are union of strata, such that
$f$ induces $g : X' \to Y'$) can be expressed also in terms of
standard operations between homotopy type. Finally, all these
constructions, still partially hypothetical (there is work on
the foundations on the drawing board!) should be able to be
paraphrased in the framework of excellent schemes, using the
machinery of non-étale tubular neighbourhoods. In both cases,
the ``fine homotopy type'' of a tame space, respectively of an
excellent scheme, is defined by passage to the limit from ``$P$-
homotopy type'' associated to finer and finer equisingular
stratifications $P$ (with non-singular strata). This ``finer''
homotopy type would necessitate the knowledge, not only of sheaves
or locally constant n-stacks, but (via a passage to the inductive

\margin{141}
